This thesis develops a vision-first training intelligence stack for football applications where a launcher is expected to target body parts and geometric zones in a closed-loop mode. Existing launcher systems in practical sport centers are typically open loop: the machine can vary speed and direction, but player state, ball trajectory, and hit events are not reconstructed in one unified 3D frame with measurable error bounds. The central objective of this work is to build, validate, and stress-test a four-camera 3D perception pipeline that can become the control backbone of an intelligent launcher.

The work spans the full \texttt{Project\_Cam} repository lifecycle. The early phase (\texttt{src/legacy}, \texttt{src/core}, \texttt{src/calibration}) established first-principles multi-view geometry and 3D rendering. The capture-hardening phase (\texttt{GARAGE\_CAMERAS}) addressed practical recording reliability and camera-device instability. The transfer phase (\texttt{garage-20260217T113109Z-3-001}) provided prior arena reconstruction logic, AprilTag mapping assets, and high-speed inference references. The integration phase (\texttt{garage\_lab\_combined}) consolidated calibration, synchronization, 3D triangulation, rendering, and quantitative GT evaluation into one reproducible workflow.

The final system includes: ChArUco intrinsics calibration, robust AprilTag extrinsics calibration, synchronized four-camera acquisition at 1280x720, per-frame 3D triangulation of ball and 17-keypoint human pose, arena-aware visualization, and protocol-based GT analysis. Ball evaluation used a 36-point static grid and dynamic clips (\texttt{ball\_slow}, \texttt{ball\_fast}, \texttt{no\_ball}). Joint evaluation used an 81-point plan (62 valid trials) for \texttt{left\_shoulder}, \texttt{right\_hip}, and \texttt{right\_knee} touch tests.

Measured results show that the stack is operational and quantitatively characterized. Intrinsics reprojection errors are sub-pixel to low-pixel. Extrinsics quality is acceptable for three cameras and weaker for one camera under deployment disturbances, emphasizing camera-stability and tag-visibility sensitivity. Ball static localization reports mean error 150.77 mm raw and 95.17 mm after axis-wise correction. Dynamic performance is stable in slow motion and degrades in fast throws due to visibility and reprojection outliers. Joint localization shows approximately 143 mm mean error, with clear per-joint bias patterns and non-negligible human placement uncertainty.

The thesis contribution is a validated perception and evaluation backbone, not a fully closed-loop launcher product. The final stage (voice command -> target semantic binding -> ballistic solver -> launcher actuation -> visual hit verification) is specified and technically feasible, but remains future work. The manuscript therefore positions the project at the transition from robust perception research to hardware-in-the-loop intelligent actuation.
